\exnumber{00}
\extitle{UI II - Ex 1 - Exploration}
\exfiles{ex00/fichier Figma, nomme Name\_FirstName\_DeliverableName\_Date\_VersionNumber}
\exauthorize{Those named in the exercise}
\makeheaderfiles

\textbf{Contexte :}

Welcome to the 2nd quest of UI!

This series of exercises is a great occasion for you to discover how libraries are made in UI, thanks to Atomic Design.

\textbf{Instructions :}

Among these 4 websites, make snapshots of these elements:

\begin{itemize}
\item 1 template
\item 1 organism
\item 1 molecule
\item 1 atom
\end{itemize}

\textbf{Websites:}

\begin{itemize}
\item Alibaba
\item Booking.com
\item Easyjet App
\item Interrail
\end{itemize}

\textbf{Livrables :}

\begin{itemize}
\item A Figma file of 4 pages
\item On each page, a screenshot of a template, an organism, a molecule and an atom
\item Déposez le fichier Figma sur Github, nommé ainsi : « Name\_FirstName\_DeliverableName\_Date\_VersionNumber ». Exemple : « Doe\_John\_MidFiPrototype\_05242024\_V1 ».
\item Ajoutez un titre dans le fichier.
\item Soignez l'aspect global du fichier. Il doit être clair, simple et lisible. Vous pouvez regarder canva.com, sans surcharger le dessin de détails.
\end{itemize}

\textbf{Ressources :}

\begin{itemize}
\item Atomic Design Methodology \url{https://atomicdesign.bradfrost.com/chapter-2/}
\item Video tutorial on Figma \url{https://www.youtube.com/watch?v=FTFaQWZBqQ8}
\item Beginning graphic design \url{https://edu.gcfglobal.org/en/beginning-graphic-design/}
\item Gareth graphic studio \url{https://www.youtube.com/c/GarethDavidStudio}
\item Example here \url{https://www.figma.com/file/fzqDkk8Zc8qJCiSMbKV0Ei/UI-II---Ex-1?node-id=0%3A1}
\end{itemize}

\textbf{Citation :}

\begin{itemize}
\item User Interface Design is a hybrid role and can bring together concepts from interaction design, visual design, information architecture and even Front End Development
\end{itemize}

